\section{Introduction}\label{sec:intro}

Most of what a group knows is worked out in conversation and stays there. Discourse seldom turns into anything of
value to the people who took part in it, and there are three reasons. A conversation leaves no object: when it ends,
nothing persists that can be used again. Where an object is made, a name for a recurring problem or a procedure
that can be invoked in a sentence, anyone can copy it, so no price forms and nothing flows back to whoever made
it. And the work is joint: one person coins a phrase, another objects to it, a third is the first to use it, and
no one of them is its owner.

A group with a private motive needs no outside payment for this work. A company that structures its own knowledge
keeps the benefit. But structured knowledge is a public good: one party's use of it does not diminish another's,
and most of its benefit falls on people other than those who made it. Markets underpay for work of this kind. The
institutions that sustained it historically, the craft guilds, paid for it by excluding outsiders
(Section~\ref{sec:prior}). Networks that have tried to pay for it directly pay on votes, on the price of a token
attached to a person, on scores that validators assign, or on counts of retrievals. Each is an opinion about
knowledge or a proxy for its use, and an opinion costs less to produce than the knowledge it concerns.

What is needed first is a unit of knowledge whose use is a fact. Such a unit can be told apart from others and
traced to an origin, and any third party can repeat and check its application without trusting whoever reports it.
What is needed second is a measure of that unit's worth that a small group cannot manufacture.

\paragraph{Objective.} This paper specifies a protocol for the first stage of a supply chain of knowledge: the
stage at which discourse is given a form that others can use. It asks two questions. How can people who take part
in making sense of something come to hold a share of what is made? And how can collaborative work accrue value to
the community that did it? We state the conditions under which serving the network pays better than gaming it, and
what remains open to an adversary when they hold.

\paragraph{Proposal.} The unit is the \emph{construction} of Fluid Construction Grammar: a pairing of a form with a
meaning and an effect. The protocol requires that applying a construction to a text be deterministic, so that a
use is a computation anyone holding the construction and the text can repeat. A set of constructions is a grammar,
and a grammar is a symbolic harness for agents: it reads what is said and makes steps available, and the agent
acts within them. A \emph{Bonfire} is a community with a grammar. Anyone who takes part in its discourse is a
member, and members whose contributions become constructions earn stake in it. A Bonfire publishes its grammar for
anyone to use and signals its grammar and what it is for. When another Bonfire meets something it cannot read, a
lookup finds a grammar that can; the constructions that did the reading are tried, and those that go on working
are adopted. A construction's worth is its \emph{breadth}: the number of independent Bonfires that have adopted it
and still hold it. Issuance is paid to originating Bonfires according to breadth, with nothing below a threshold,
so that what a few parties agree among themselves earns nothing. A share of what a construction earns passes to
the parents it declares. \emph{Nodes} supply the computation, and the stake they pledge to a Bonfire sets a ceiling
on what it can be paid, which only breadth can fill. What cannot be computed is whether several Bonfires are one
party. That is answered by stake that can be lost, and under the assumptions of Section~\ref{sec:security} posing as
many does not pay so long as the ceiling rate
stays below the cost of capital plus the expected loss from being caught.

\paragraph{Scope.} This paper is a specification with an argument. It supersedes the first version
\citep{cortes2025value}. Constructions, their teaching and a tally of their uses inside a single group are built and
described in a companion paper \citep{vera2026constructions}; the network of Bonfires and nodes is proposed here
(Appendix~\ref{app:built}), with the questions it leaves open in Section~\ref{sec:open}. The protocol is
indifferent to whether discourse is produced by people or by agents: the distinction cannot be checked, and the
protocol asks only whether a construction does checkable work for communities other than its own. Supply and
issuance schedules are specified in a separate token model document; here we refer only to \emph{the epoch
budget}, the amount divided in one period.
